% Inter-rater agreement appendix

\subsection{Design}
Inter-rater agreement (IAA) was assessed on a stratified subsample of 1{,}800 samples drawn from three evaluation batches (50 random samples per dataset per batch).
Each sample received three independent DeepSeek (\texttt{deepseek-chat}) annotations using the same Pass~1 and Pass~2 templates as the full-scale evaluation. Krippendorff's ordinal agreement coefficient $\alpha_{\mathrm{ord}}$ was computed on three complete ratings. Pass~1 covered all 1{,}800 rows. Pass~2 effect dimensions used $n{=}1{,}147$ samples with three complete severity ratings.
All three annotations used the same judge model, so $\alpha_{\mathrm{ord}}$ reflects scoring variation under repeated sampling, not agreement between human raters.
% Low Pass~1 agreement on behavioral resistance ($\alpha_{\mathrm{ord}}{=}0.38$) suggests that category is ambiguous even under identical rubrics, which may contribute to Pass~1 false-positives (Section~\ref{sec:results}).

\subsection{Pass~1 Agreement}
Table~\ref{tab:iaa_pass1} reports ordinal $\alpha_{\mathrm{ord}}$ for Pass~1 screening categories.
Steerability, logical coherence, and factual integrity reached moderate-to-substantial agreement ($\alpha_{\mathrm{ord}} \ge 0.75$).
Behavioral resistance was lower ($\alpha_{\mathrm{ord}}{=}0.38$), reflecting that category is ambiguous even under identical rubrics, which may contribute to Pass~1 false-positives.

\begin{table}[htbp]
  \centering\small
  \caption{Pass~1 inter-rater agreement (ordinal $\alpha_{\mathrm{ord}}$, three raters).}
  \label{tab:iaa_pass1}
  \begin{tabular}{@{}lr@{}}
    \toprule
    Category & $\alpha_{\mathrm{ord}}$ \\
    \midrule
    Steerability & 0.80 \\
    Logical coherence & 0.81 \\
    Factual integrity & 0.75 \\
    Output shape & 0.64 \\
    Tonal change & 0.64 \\
    Behavioral resistance & 0.38 \\
    \bottomrule
  \end{tabular}
\end{table}

\subsection{Pass~2 Agreement}
Table~\ref{tab:iaa_pass2} lists Pass~2 side-effect dimensions sorted by agreement.
Dimensions above the 0.65 threshold used for primary claims selection include repetition looping, inverse logit polarity, factual reversal, option-label mismatch, token inflation, factual drift, steering asymmetry, and label-content contradiction.
Logit-text decoupling ($\alpha_{\mathrm{ord}}{=}0.39$), tonal rigidity ($0.49$), and hedging escalation ($0.50$) fell below this threshold and are treated as exploratory findings despite frequent BH significance in the full-scale evaluation.

\begin{table}[htbp]
  \centering\small
  \caption{Pass~2 inter-rater agreement (ordinal $\alpha_{\mathrm{ord}}$, $n{=}1{,}147$ samples, three raters).}
  \label{tab:iaa_pass2}
  \begin{tabular}{@{}lr@{}}
    \toprule
    Dimension & $\alpha_{\mathrm{ord}}$ \\
    \midrule
    Inverse logit polarity & 0.87 \\
    Repetition looping & 0.84 \\
    Factual reversal & 0.77 \\
    Option-label mismatch & 0.72 \\
    Token inflation & 0.71 \\
    Factual drift & 0.69 \\
    Steering asymmetry & 0.69 \\
    Label-content contradiction & 0.68 \\
    Hallucinated specificity & 0.67 \\
    Sensitivity refusal & 0.66 \\
    Persona amplification & 0.60 \\
    Implicit assumption injection & 0.54 \\
    Safety override resistance & 0.53 \\
    Token compression & 0.51 \\
    Monotonic factual scaling & 0.50 \\
    Hedging escalation & 0.50 \\
    Tonal rigidity & 0.49 \\
    Neutralization rephrasing & 0.46 \\
    Logit-text decoupling & 0.39 \\
    \bottomrule
  \end{tabular}
\end{table}

As a supplementary check, we compared judge-flagged inverse logit polarity against polarity inferred directly from target-token logits on all 16{,}015 Pass~2 scored rows (precision 1.0, recall 0.99).
This objective oracle aligns with the high IAA on this dimension.

\subsection{Overall vs.\ Per-\texorpdfstring{$\alpha$}{alpha} Validation}\label{per alpha result}
Section~\ref{permultiplierresult} summarizes automated per-multiplier proxies on the full-scale evaluation, while judge-based per-$\alpha$ validation statistics are reported here.
Full-scale Pass~2 uses overall cross-strength severity scores.
To assess whether overall scores track decomposed per-strength ratings, we re-annotated a stratified subsample of 400 Pass~1-flagged inputs with three judges using the per-$\alpha$ judge template, covering five dimensions: steering asymmetry, token inflation, factual reversal, inverse logit polarity, and repetition looping.
Of 400 inputs, 262 samples yielded comparable overall Pass~2 scores from the full-scale evaluation and complete per-$\alpha$ triple ratings.
The remainder lacked parseable per-$\alpha$ judge output or mismatched run identifiers.

Table~\ref{tab:overall_per_alpha} reports Pearson correlations between overall Pass~2 scores and max(per-$\alpha$) severity on the comparable subset; max(per-$\alpha$) denotes each sample's peak per-strength severity, averaged across samples.
Repetition looping tracked max(per-$\alpha$) closely ($r{=}0.96$).
Token inflation ($r{=}0.73$) and factual reversal ($r{=}0.59$-$0.65$) showed moderate-to-strong alignment, and steering asymmetry was weaker ($r{=}0.44$), with overall scores systematically higher than max(per-$\alpha$) severity (2.51 vs.\ 1.98).
These correlations support using automated per-multiplier proxies and the subsample to interpret overall Pass~2 scores, while acknowledging dimension-dependent bias.

\begin{table}[htbp]
  \centering\small
  \caption{Overall vs.\ max(per-$\alpha$) severity correlations ($n{=}262$ comparable samples, three raters).}
  \label{tab:overall_per_alpha}
  \resizebox{\linewidth}{!}{%
  \begin{tabular}{@{}lrrrr@{}}
    \toprule
    Dimension & $r$ (Pearson) & MAE & Overall mean & max($\alpha$) mean \\
    \midrule
    Repetition looping & 0.96 & 0.13 & 0.87 & 0.89 \\
    Token inflation & 0.73 & 1.10 & 1.69 & 2.31 \\
    Factual reversal & 0.59 & 0.90 & 1.34 & 1.10 \\
    Inverse logit polarity & 0.53 & 0.71 & 0.90 & 0.68 \\
    Steering asymmetry & 0.44 & 1.30 & 2.51 & 1.98 \\
    \bottomrule
  \end{tabular}%
  }
\end{table}

\subsection{Parser and Data Notes}
A parser edge case affected 34 Pass~2 rows where JSON evidence listed non-zero inverse logit polarity severity but exported CSV columns read zero.
These rows are excluded from the six incomplete Pass~2 cases documented in Section~\ref{Conditional Evaluation Sample}.
The issue was fixed in \texttt{pipeline/judge\_parser.py} for future runs.

\subsection{Pass~1\texorpdfstring{$\to$}{->}Pass~2 Concordance}
Pass~1$\to$Pass~2 concordance is defined in Section~\ref{Two-Pass Evaluation Protocol}.
Among Pass~1-flagged samples, 89.7\% had at least one non-zero Pass~2 severity; 1{,}649 (10.3\%) were flagged in Pass~1 but received all-zero Pass~2 scores.
Pass~1 steerability predicts Pass~2 magnitude: \textsc{Moderate}/\textsc{Strong}/\textsc{Anti-Steered} labels yielded $\ge 98\%$ non-zero Pass~2 rates and mean max severity $\ge 3.6$, whereas \textsc{Weak} flagged samples had 84.0\% non-zero and mean max severity 2.38 (mean per-sample maximum across the 19 Pass~2 dimensions).
Figure~\ref{fig:steerability_asymmetry_box} shows the same pattern for steering asymmetry.
Category-to-dimension overlap was highest for behavioral resistance$\to$steering asymmetry (77.8\% non-zero when resistance flagged, mean severity 2.42) and factual integrity$\to$factual reversal (75.1\%, mean 3.30), and lowest for behavioral resistance$\to$logit-text decoupling (3.3\% non-zero) (Figure~\ref{fig:pass1_pass2_concordance}).
Samples with more co-flagged Pass~1 categories showed higher Pass~2 severity ($\ge 3$ categories: $\ge 99\%$ non-zero).

\begin{figure}[htbp]
    \centering
    \includegraphics[width=\textwidth]{pass1_pass2_concordance_heatmap.png}
    \caption{Pass~1$\to$Pass~2 concordance: percentage of category-flagged samples with non-zero Pass~2 severity on each nested dimension (Section~\ref{sec:methodology}).}
    \label{fig:pass1_pass2_concordance}
\end{figure}

\begin{figure}[htbp]
    \centering
    \includegraphics[width=0.85\textwidth]{pass1_steerability_vs_asymmetry_boxplot.png}
    \caption{Boxplots of Pass-2 steering asymmetry severity (y-axis; ordinal 0-5) by Pass-1 steerability label (x-axis) among Pass-1-flagged samples ($N{=}16{,}015$). Pass-1 labels follow Section~\ref{sec:methodology}; \textsc{Unmeasurable} denotes samples whose baseline output precludes a valid steerability comparison across strengths.}
    \label{fig:steerability_asymmetry_box}
\end{figure}

\subsection{Supplementary Severity Figures}
Figure~\ref{fig:pass1_flags} shows Pass~1 any-category flag rate by behavioral dataset on the full corpus ($N{=}28{,}807$).
Figure~\ref{fig:heatmap} shows mean Pass~2 severity across the full dataset$\times$dimension grid among Pass~1-flagged samples ($N{=}16{,}015$).
Figure~\ref{fig:top_dims} ranks dimensions by average mean severity across datasets.
The four primary claim dimensions appear among the highest-ranked effects (Table~\ref{tab:dimension_claims}).
Figure~\ref{fig:score_dist} shows Pass~2 severity distributions for the eight highest-mean dimensions among Pass~2-evaluated samples ($N{=}16{,}015$).

\begin{figure}[htbp]
    \centering
    \includegraphics[width=0.7\textwidth]{pass1_flag_rate_by_dataset.png}
    \caption{Pass-1 any-category flag rate (horizontal axis, \%) by behavioral dataset (vertical axis). Each bar is one of 36 datasets; $N{=}28{,}807$ samples total.}
    \label{fig:pass1_flags}
\end{figure}

\begin{figure}[htbp]
    \centering
    \includegraphics[width=\textwidth]{severity_heatmap_dataset_dimension.png}
    \caption{Mean Pass-2 severity (color scale, 0-5) by behavioral dataset (rows) and side-effect dimension (columns) among Pass-1-flagged samples.}
    \label{fig:heatmap}
\end{figure}

\begin{figure}[htbp]
    \centering
    \includegraphics[width=0.85\textwidth]{top_dimensions_bar.png}
    \caption{Average mean Pass~2 severity across 36 datasets (horizontal axis) by side-effect dimension (vertical axis).}
    \label{fig:top_dims}
\end{figure}

\begin{figure}[htbp]
    \centering
    \includegraphics[width=\textwidth]{score_distribution_by_dimension.png}
    \caption{Pass~2 severity score distributions for the eight highest-mean dimensions among Pass~2-evaluated samples ($N{=}16{,}015$). Each panel plots sample count (y-axis) vs.\ ordinal severity 0-5 (x-axis).}
    \label{fig:score_dist}
\end{figure}

\subsection{Perplexity vs.\ Side-Effect Severity}
Figure~\ref{fig:perplexity_side_effects} supplements the corpus-mean perplexity curves in Section~\ref{sec:results} (Figure~\ref{fig:perplexity_strength}).
Per-sample max $|\Delta\mathrm{perplexity}|$ from the $\alpha{=}0$ baseline correlates weakly with overall Pass~2 severity on the four primary-claim dimensions (Figure~\ref{fig:perplexity_side_effects}).

\begin{figure}[htbp]
    \centering
    \includegraphics[width=\textwidth]{perplexity_vs_side_effects_scatter.png}
    \caption{Hexbin density of overall Pass~2 severity (y-axis, ordinal 0-5) vs.\ per-sample max $|\Delta\mathrm{perplexity}|$ from the $\alpha{=}0$ baseline (x-axis) on the Pass~2-evaluated subset ($N{=}16{,}015$). Panels show the four primary-claim dimensions. Color encodes sample count per bin. Annotated $\rho$ is Pearson correlation.}
    \label{fig:perplexity_side_effects}
\end{figure}
